% Generated by `taskgen course examples` from number_theory.primitive_root (seed 4).
% Do not edit: change examples.yaml and regenerate.
\begin{example}[поиск наименьшего первообразного корня, $n = p$]\label{ex:primroot-search}
Найти наименьший первообразный корень по модулю $n = 193$ или доказать, что первообразных корней по модулю $193$ нет.

\par\noindent\textit{Решение.} $n = 193$ --- простое, поэтому первообразные корни существуют. $\varphi(n) = 192 = 2^{6} \cdot 3$; $g$ ($\gcd(g, n) = 1$) --- первообразный корень $\iff g^{\varphi(n)/q} \not\equiv 1$ для всех простых $q \mid \varphi(n)$, т.\,е. для показателей $192/2 = 96$, $192/3 = 64$. Перебираем $g = 2, 3, \dots$ (квадраты $h^2$ пропускаем: $(h^2)^{\varphi(n)/2} = h^{\varphi(n)} \equiv 1$):\[\begin{array}{c|cc|l}g & g^{96} & g^{64} & \\ \hline 2 & 1 &  & \text{нет} \\ 3 & 1 &  & \text{нет} \\ 4 &  &  & = 2^2 \\ 5 & 192 & 84 & \text{корень} \\ \end{array}\](вычисление прекращаем на первом показателе, дающем $1$.)\[g = 5; \quad \text{всего первообразных корней } \varphi(\varphi(n)) = \varphi(192) = 64.\]
\end{example}
